\documentclass[../ap_2012.tex]{subfiles}

\graphicspath{{./image/}}

\begin{document}

\setcounter{chapter}{1}
\chapter{線形代数:}
\section{}
固有値方程式より、
\begin{equation}\begin{aligned}[b]
    0 &= \abs{A-\lambda I}\\
    &= \begin{vmatrix}
        -\lambda & -1 & -1\\
        -1 & -\lambda & -1\\
        -1 & -1 & \alpha-\lambda
    \end{vmatrix}
    &=\lambda^2(\alpha-\lambda)-2+2\lambda-(\alpha-\lambda)\\
    &=(\lambda^2-1)(\alpha-\lambda)+2(\lambda-1)\\
    &=(\lambda-1)(-\lambda^2+(\alpha-1)\lambda+\alpha+2)\\
    \lambda &= 1,\,\frac{\alpha-1\pm\sqrt{(\alpha+1)^2+8}}{2}
\end{aligned}\end{equation}
いま、\(\alpha=5/2\)より、
\begin{equation}\begin{aligned}[b]
    \lambda &= 1, \frac{3}{4}\pm\frac{9}{4}\\
    &=1,\,3,\,-\frac{3}{2}
\end{aligned}\end{equation}

\(\lambda=1\)に対応する固有ベクトルは
\begin{equation}\begin{aligned}[b]
    \begin{pmatrix}
        -1 & -1 & -1\\
        -1 & -1 & -1\\
        -1 & -1 & 3/2
    \end{pmatrix}\begin{pmatrix}
        a\\b\\c
    \end{pmatrix}
    =0
\end{aligned}\end{equation}
より、\(v_1=(1,-1,0)^\mathsf{T}/\sqrt{2}\)

\(\lambda=3\)に対応する固有ベクトルは
\begin{equation}\begin{aligned}[b]
    \begin{pmatrix}
        -3 & -1 & -1\\
        -1 & -3 & -1\\
        -1 & -1 & -1/2
    \end{pmatrix}\begin{pmatrix}
        a\\b\\c
    \end{pmatrix}
    =0
\end{aligned}\end{equation}
より、\(v_2=(1,1,-2)^\mathsf{T}/\sqrt{6}\)

\(\lambda=-3/2\)に対応する固有ベクトルは
\begin{equation}\begin{aligned}[b]
    \begin{pmatrix}
        3/2 & -1 & -1\\
        -1 & 3/2 & -1\\
        -1 & -1 & 4
    \end{pmatrix}\begin{pmatrix}
        a\\b\\c
    \end{pmatrix}
    =0
\end{aligned}\end{equation}
より、\(v_3=(2,2,1)^\mathsf{T}/3\)

\section{}
\subsection{}
前問より
\begin{equation}\begin{aligned}[b]
    \lambda &= 1,\,\frac{\alpha-1\pm\sqrt{(\alpha+1)^2+8}}{2}
\end{aligned}\end{equation}
であり、
\begin{equation}\begin{aligned}[b]
    (\alpha+1)^2+8 > 0
\end{aligned}\end{equation}
より固有値は全て実数。
\subsection{}
\begin{equation}\begin{aligned}[b]
    \lambda &= 1,\,\frac{\alpha-1\pm\sqrt{(\alpha+1)^2+8}}{2}\\
    &\to 1,\,\frac{\alpha-1}{2}\pm\frac{\alpha+1}{2}\qquad(\alpha\to\infty)\\
    &= 1,\,-1\,\alpha
\end{aligned}\end{equation}
より
\begin{equation}\begin{aligned}[b]
    \lim_{\alpha\to\infty}\frac{\max(\lambda)}{\alpha} = 1
\end{aligned}\end{equation}

\subsection{}
\begin{equation}\begin{aligned}[b]
    \lim_{\alpha\to\infty}\frac{\min(\lambda)}{\alpha} = -1
\end{aligned}\end{equation}

\section{}
\(\alpha=0\)のときの固有値は
\(\lambda=1,-2\)で、\(\lambda=1\)は縮退している。
\(\lambda=1\)に対応する固有ベクトルは
\begin{equation}\begin{aligned}[b]
    \begin{pmatrix}
        -1 & -1 & -1\\
        -1 & -1 & -1\\
        -1 & -1 & -1
    \end{pmatrix}\begin{pmatrix}
        a\\b\\c
    \end{pmatrix}
    =0
\end{aligned}\end{equation}
より、\(v_1=(1,-1,0)^\mathsf{T}/\sqrt{2},v_2=(0,1,-1)^\mathsf{T}/\sqrt{2}\)

\(\lambda=-2\)に対応する固有ベクトルは
\begin{equation}\begin{aligned}[b]
    \begin{pmatrix}
        2 & -1 & -1\\
        -1 & 2 & -1\\
        -1 & -1 & 2
    \end{pmatrix}\begin{pmatrix}
        a\\b\\c
    \end{pmatrix}
    =0
\end{aligned}\end{equation}
より、\(v_3=(1,1,1)^\mathsf{T}/\sqrt{3}\)
となる。
いま、\(A\)は対称行列なので、直交行列\(P=(v_1,v_2,v_3)\)を使って
\begin{equation}\begin{aligned}[b]
    \begin{pmatrix}
        1 & 0 & 0\\
        0 & 1 & 0\\
        0 & 0 & -2
    \end{pmatrix}
    =P^\mathsf{T}AP
\end{aligned}\end{equation}
のように対角化できる。

よって
\begin{equation}\begin{aligned}[b]
    A^n &= P\qty(P^\mathsf{T}AP)^n P^\mathsf{T}
    =\begin{pmatrix}
        v_1 & v_2 & v_3
    \end{pmatrix}\begin{pmatrix}
        1 & 0 & 0\\
        0 & 1 & 0\\
        0 & 0 & (-2)^n
    \end{pmatrix}\begin{pmatrix}
        v_1^\mathsf{T}\\
        v_2^\mathsf{T}\\
        v_3^\mathsf{T}\\
    \end{pmatrix}\\
    &=\begin{pmatrix}
        v_1 & v_2 & (-2)^nv_3
    \end{pmatrix}\begin{pmatrix}
        v_1^\mathsf{T}\\
        v_2^\mathsf{T}\\
        v_3^\mathsf{T}\\
    \end{pmatrix}
    = v_1v_1^\mathsf{T} + v_2v_2^\mathsf{T} + (-2)^nv_3v_3^\mathsf{T}\\
    &= I + ((-2)^n-1)v_3v_3^\mathsf{T}
    = \begin{pmatrix}
        1 & 0 & 0\\
        0 & 1 & 0\\
        0 & 0 & 0
    \end{pmatrix}
    +\frac{(-2)^n-1}{3}\begin{pmatrix}
        1 & 1 & 1\\
        1 & 1 & 1\\
        1 & 1 & 1
    \end{pmatrix}\\
    &=\frac{1}{3}\begin{pmatrix}
        (-2)^n+2 & (-2)^n-1 & (-2)^n-1\\
        (-2)^n-1 & (-2)^n+2 & (-2)^n-1\\
        (-2)^n-1 & (-2)^n-1 & (-2)^n+2
    \end{pmatrix}
\end{aligned}\end{equation}

\end{document}
